\chapter{Collaboration and Next Step}

\section{Where a cofounder has immediate leverage}

The value of a cofounder at this stage is not just vision alignment.
It is making the wedge commercially real.

There are four areas where the contribution matters most right now.

\subsection{Shaping the wedge under market pressure}

The wedge is a hypothesis until a customer pays.
A cofounder who is willing to have direct conversations helps pressure-test:

\begin{itemize}
  \item which workflows are most painful,
  \item which segments move fastest,
  \item which version of the story gets meetings,
  \item which version of the pitch reduces objections.
\end{itemize}

\subsection{Identifying pilot customers}

The most valuable short-term contribution is identifying two to three credible pilot opportunities where:

\begin{itemize}
  \item AI is already live in a real workflow,
  \item the decision has consequences,
  \item the buyer can move within weeks,
  \item the workflow can be scoped tightly.
\end{itemize}

This matters more than broad lead generation.
Quality of the first conversations is more important than volume.

\subsection{Refining GTM and positioning}

The practical GTM questions need answers:

\begin{itemize}
  \item Who do we call first?
  \item Which segment lands fastest?
  \item What wording creates immediate recognition of the pain?
  \item What pilot package is easiest to approve?
  \item What outcome do we promise and how do we measure it?
\end{itemize}

These are not strategic questions to revisit quarterly.
These are operational questions to answer in weeks and refine continuously.

\subsection{Compressing the learning cycle}

We need fast feedback loops from real conversations.
A cofounder accelerates this because:

\begin{itemize}
  \item two people can have more conversations simultaneously,
  \item two perspectives on the same conversation improve interpretation,
  \item faster iteration on the pitch means faster clarity on what works.
\end{itemize}

\section{Roles}

\begin{center}
\renewcommand{\arraystretch}{1.6}
\begin{tabularx}{0.92\textwidth}{@{}lXX@{}}
\toprule
\textbf{Area} & \textbf{Founder} & \textbf{Cofounder} \\
\midrule
Product and architecture & Primary & Advisory \\
Pilot design and scoping & Primary & Supporting \\
Outreach and pipeline & Shared & Shared \\
Positioning and messaging & Shared & Shared \\
Pilot customer identification & Supporting & Primary \\
Network activation & Supporting & Primary \\
\bottomrule
\end{tabularx}
\end{center}

\section{The next step}

The next step is narrow and concrete.

\begin{conceptbox}{Next step}
Identify two to three pilot opportunities and validate the wedge in real workflows.
\end{conceptbox}

Specifically, we need to confirm:

\begin{enumerate}
  \item \textbf{Who we call first.} Which teams in which segments are most likely to feel the pain right now?
  \item \textbf{What we sell first.} Which workflow framing gets the fastest yes?
  \item \textbf{Why they will pay.} Which version of the risk-reduction story closes the conversation?
\end{enumerate}

If we can answer those three questions with real customer conversations and one or two paid pilots, the rest of the commercial path becomes much clearer.

\section{What this plan is not claiming}

This plan is not claiming that every assumption here is correct.
Some segments may move slower than expected.
Some framing may need adjustment.
The pilot timeline may vary.

That is fine.
The purpose of this plan is to create enough clarity to take action.

Once we are in real conversations, we update based on what we learn.
That is how an early market entry actually works.

\section{Closing}

The purpose of this document was simple:

\begin{itemize}
  \item Explain where Connector enters.
  \item Explain what we sell.
  \item Explain why someone pays.
  \item Explain what a pilot looks like.
  \item Explain how we go to market.
  \item Explain where a cofounder plugs in.
\end{itemize}

If those six things are clear after reading this, the document did its job.

\begin{insightbox}{One-line summary}
\centering
\large ``Pick one AI workflow where decisions matter. We make it traceable, explainable, and provable. Then we expand.''
\end{insightbox}

That is the whole strategy.
Everything else is execution.
